% ABSRACT PAGE
\chapter*{Abstract}
\begin{spacing}{1.5}
	
The Abstract serves as a condensed version of your entire project. According to the syllabus, it must reflect the process of identifying a need and converting it into a problem statement. It acts as a standalone summary that allows a reader to understand the essence of your Web-Based ML project without reading the full document.

Abstract Length Guidelines
For a Mini Project of this scale, the abstract should typically be between 250 to 300 words. It should fit on a single page and be written as a single, continuous paragraph or a few short, dense paragraphs.

Detailed Guidelines for Writing the Abstract
To meet the assessment criteria for clarity and societal impact, follow this structural flow:

1. Context and Motivation (2–3 sentences)
Start with the domain (Web + ML).
State the societal or research need identified during your survey.

Example: "With the rise of digital health, there is a significant need for accessible diagnostic tools for rural populations..."

2. Problem Statement (1–2 sentences)

Clearly define the technical problem you are solving.

Mention the specific gap identified in existing systems.

3. Proposed Solution and Methodology (3–4 sentences)

Identify the specific ML algorithms and web technologies used.

Briefly mention the data source and the "best possible solution" selected.

Example: "This project proposes a Flask-based web application integrated with a Random Forest model trained on the XYZ dataset..."

4. Results and Validation (2–3 sentences)

Highlight the "full functioning of the working model".

Mention key performance metrics (accuracy, efficiency) obtained through simulations or testing.

5. Impact and Conclusion (1–2 sentences)

State how the project achieves "sustainable development" or maps to an SDG.

Summarize the cost-effectiveness and societal value of the final product.

Abstract "Checklist" for Students
No Citations: Do not include references (e.g., [1]) in the abstract.
No Figures/Tables: Keep it strictly text-based.

Past Tense: Since the abstract summarizes work already done, use the past tense (e.g., "The model was trained..." or "The survey was conducted...").

Keywords: Include 4–6 keywords at the bottom (e.g., Machine Learning, Web Development, SDG 3, Predictive Modeling).
\newline	\\
\textbf{Keywords:} AI,  Machine Learning, (include atleast five keywords) 	
	
\end{spacing}